\section{Unbinned background subtraction by negative-weight injection}
\label{app:negweight}

Background subtraction in \S\ref{sec:method-algo} is the only step that uses
analysis binning before final histogramming. Each data event in reconstructed
bin $b$ receives the purity weight
$w_{\mathrm{pur}}(b)=\max(0,\,N_{\mathrm{data}}(b)-N_{\mathrm{bkg}}(b))/
N_{\mathrm{data}}(b)$. Injecting background simulation on the measured side
with negative weight gives a fully unbinned correction with the same
bin-integrated estimand. Its agreement with the binned correction is quantified
on the 2D measurement below. It is the continuous version of the purity weight
recommended for nontrivial backgrounds in Ref.~\cite{unbinnedguide}, and the
natural treatment for the larger backgrounds of the extended-fiducial analysis.

\paragraph{Scope of the numbers.}
No number here is a publication uncertainty product; the analytic derivation,
estimand identity and displacement argument require none of them. Every result
in this note uses the binned purity weight of \S\ref{sec:method-algo}.
Negative-weight injection serves as an equivalence study and the natural
extended-fiducial treatment; this appendix establishes no supported production
path and changes no default. The two
covariance ratios are agreement diagnostics between two realizations of one
subtraction, not uncertainty products. Their provenance and the two limits
attached to the production numbers are in App.~\ref{app:provenance}
(\texttt{negweight-hpss}, \texttt{negweight-toy}).

\subsection{Estimator}
\label{sec:negweight-estimator}

OmniFold step 1 trains a classifier to separate simulated reconstruction (class~0, weighted density
$S$) from the measured sample (class~1, weighted density $\rho_1$), then
reweights simulation by the learned density ratio $r(x)=p(x)/(1-p(x))$. Its population optimum is
$r^{\star}(x)=\rho_1(x)/S(x)$. Background treatment is entirely a choice of $\rho_1$:
\begin{itemize}
  \item binned purity assigns every data event in reconstructed analysis bin
    $b$ the same scalar. For $x\in b$,
    \begin{equation}
      \label{eq:negweight-purity}
      \rho_1^{\mathrm{pur}}(x)
      \;=\; D(x)\,w_{\mathrm{pur}}(b)
      \;=\; D(x)\,\max\!\Big(0,\;1-\frac{B_{b}}{D_{b}}\Big),
    \end{equation}
    where $D$ is the data density, $B$ the POT-scaled background density,
    and $D_{b},B_{b}$ their bin masses;
  \item negative-weight injection appends reconstructed background-simulation
    events to the measured side with weight
    $-w_{\mathrm{bkg}}\,\times\,(\text{POT scale})$ and keeps data weights at
    $+1$, giving
    \begin{equation}
      \label{eq:negweight-continuous}
      \rho_1^{\mathrm{neg}}(x) \;=\; D(x)-B(x)
    \end{equation}
    pointwise, without analysis binning.
\end{itemize}
Equations~\eqref{eq:negweight-purity} and \eqref{eq:negweight-continuous}
differ as functions of $x$. A bin-constant rescaling reproduces $D(x)-B(x)$
inside bin $b$ only when $B/D$ is constant there. Otherwise,
$D(x)\,(1-B_{b}/D_{b})$ and $D(x)-B(x)$ share a bin integral but differ
within the bin. The constructions agree \emph{exactly} on bin masses under the
piecewise-constant sieve of \S\ref{sec:negweight-binned}; in general they differ
by within-bin shape terms. Both target the \emph{bin-integrated} estimand
$r^{\star}(b)=(D_{b}-B_{b})/S_{b}$; they do not share the continuous estimand
$r^{\star}(x)=(D(x)-B(x))/S(x)$ (App.~\ref{app:history}).
\S\ref{sec:negweight-2d} and \S\ref{sec:negweight-uq} measure the effect
of that difference. The injected background and data occupy the same
reconstructed fiducial window, giving identical support.

\paragraph{The floor.} Where $D-B<0$, binned purity floors the correction at
zero through $\max(0,\cdot)$. The unbinned classifier's bounded output
$p\in(0,1)$ forces $r=p/(1-p)\ge 0$. Non-negativity is structural, while the value at the floor follows from the
bounded output and ordinary regularization rather than from a minimizer of
the fit (\S\ref{sec:negweight-binned}). Locally negative
class-1 mass makes the fit ill-posed and calls for a positive-target
prescription. We apply the ``Stay Positive'' refinement of
Ref.~\cite{Nachman:2025staypositive}, then explicitly zero-clip the learned
refinement factor before step-1 training. Agreement between plain injection
and clipped refinement depends on the classifier.

For the histogrammed gradient-boosted density estimator used in \emph{this
comparison}, they agree to the reported median $\nwMedianBin$ at this
measurement's sub-percent genuine-background scale. The \texttt{hist} backend
differs from both the production central value's \texttt{exact} and the
covariance ensemble's \texttt{lgbm} (\S\ref{sec:method}). This validation therefore establishes the equivalence on a
third backend, and its transfer to the production configuration is an assumption
rather than a demonstration.
Raw signed weights destabilize an \emph{exact} (tree-boosted) classifier,
inflating a full low-$p_T$ cross-section row by up to a factor
${\sim}\nwSpRawMax$ where local background is large. Clipped Stay-Positive
refinement removes this pathology and recovers binned-purity agreement to a
median $\nwSpMedian$ ($\nwSpWithin$ bins within \SI{5}{\percent}).
This clipped refinement is
therefore what makes the unbinned subtraction robust across estimators, and it
matters most for the larger, more structured backgrounds of the
extended-fiducial analysis.

\subsection{Binned limit: exact reduction to the purity correction}
\label{sec:negweight-binned}

A step-1 classifier restricted to functions piecewise-constant on reconstructed
analysis bins $\{b\}$ gives the population-level binned limit of the
OmniFold/Richardson--Lucy correspondence
in Refs.~\cite{Andreassen:2019cjw,Andreassen:2021zzk}. Its output is a single
$p$ per bin. With common support and unnormalized class weights, step~1
weighted cross-entropy separates into per-bin objectives
\begin{equation}
  \label{eq:negweight-Lb}
  L_b(p) = -W_1\,\ln p - W_0\,\ln(1-p),
\end{equation}
where $W_0=S(b)>0$ is the integral of simulated class-0 density $S$ in bin $b$,
and $W_1$ is the signed class-1 weight. Injecting background simulation at
negative weight while keeping data at $+1$ gives $W_1=D(b)-B(b)$.
Here $D(b)$ and $B(b)$ are the bin masses of $D$ and $B$, equivalently
$N_{\mathrm{data}}(b)$ and $N_{\mathrm{bkg}}(b)$ in
\S\ref{sec:negweight-estimator}.

For $D(b)>B(b)$, objective~\eqref{eq:negweight-Lb} is strictly convex on
$(0,1)$ with unique interior minimizer $p=W_1/(W_0+W_1)$, giving
\begin{equation}
  \label{eq:negweight-rb}
  r(b) = \frac{p}{1-p} = \frac{W_1}{W_0} = \frac{D(b)-B(b)}{S(b)}.
\end{equation}
Purity gives exactly this reweight under the same sieve, since its class-1
weight $W_1=D(b)\,w_{\mathrm{pur}}(b)=\max(0,\,D(b)-B(b))$ equals $D(b)-B(b)$
on this domain. At $D(b)=B(b)$, the infimum lies at $p\to0^+$ and
$r(b)\to0$, again giving purity when the boundary is admitted.
The treatments therefore agree bin by bin at the population optimum under the
common sieve. Combined with the binned OmniFold/Richardson--Lucy correspondence,
they give the traditional subtract-background-then-unfold update.
Finite, regularized classifiers approximate this identity; their residual
differences are tested below.

The reduction fails where $D(b)<B(b)$. With $W_1<0$,
$L_b(p)=|W_1|\,\ln p-W_0\,\ln(1-p)$ tends to $-\infty$ as $p\to0^+$
and is unbounded below, with no minimizer. A minimizing sequence gives $r(b)\to0$.
The $\max(0,\cdot)$ floor is a prescription, not a minimizer's value.
A finite regularized classifier's behavior depends on the estimator.
Stay-Positive remains well behaved but can
retain negative refined weights where the target density is genuinely negative.
The
zero clip used here is therefore an additional analysis prescription chosen to
match the physical purity floor.
Ref.~\cite{Nachman:2025staypositive} documents the same reweighting-loss
divergence for negative targets; signed binary cross-entropy has an analogous
pathology for sWeights~\cite{GlazierTyson:2024}.

The binned $D-B$ update, negative-weight classification of the measured side,
and binned OmniFold/Richardson--Lucy correspondence already appear together in
Ref.~\cite{Andreassen:2021zzk}. Equation~\eqref{eq:negweight-rb} states their
immediate binwise corollary and its precise validity domain for this analysis.
It is a consistency check, with no separate method claim.

\subsection{Closed-form toy}
\label{sec:negweight-toy}

A one-dimensional toy uses signal simulation $\mathcal{N}(0,1)$, equally
normalized data signal $\mathcal{N}(0.3,1.1)$, and background
$\mathcal{N}(2,0.7)$ contributing \SI{16.7}{\percent} of pseudo-data.
The known target $(D-B)/S=\phi(x;0.3,1.1)/\phi(x;0,1)$ is smooth and positive.
A signed-weight gradient-boosted classifier recovers it with median
\SI{\nwToyNeg}{\percent} relative error, matching binned purity's
\SI{\nwToyPur}{\percent} and reproducing it bin-for-bin on the same problem.
Omitting subtraction leaves \SI{\nwToyBias}{\percent} residual-background bias.
Results are stable across classifier seeds ($\pm\SI{\nwToySeed}{\percent}$),
with no numerical instability from negative weights. Moving background into
the sparse signal tail and deliberately over-subtracting makes the learned
reweight floor cleanly at zero where $B>D$, as anticipated.

\subsection{Comparison on the 2D measurement}
\label{sec:negweight-2d}

With identical unfolding settings and a common classifier seed, the ME~FHC
comparison isolates background treatment (Table~\ref{tab:negweight}).
Total cross sections agree to \SI{\nwPctTot}{\percent}, an order of magnitude
inside the ${\sim}\SI{3}{\percent}$ purity correction. The per-bin ratio has
median $\nwMedianBin$ and \SI{\nwRmsBin}{\percent} RMS scatter across the plane
(Fig.~\ref{fig:negweight-ratio}). Only one low-$p_T$, high-background edge bin
differs by more than a few percent (\SI{\nwWorstBin}{\percent}); background
concentrates there and subtraction is largest. Unbinned subtraction reproduces
the adopted binned correction wherever it is well determined. Departures occur
only where the binned $\max(0,\cdot)$ and structural floor of
\S\ref{sec:negweight-estimator} govern the result.

\begin{table}[t]
  \centering
  \caption{Binned purity versus unbinned negative-weight background
    subtraction on the 2D ME~FHC measurement (common estimator and seed).
    The per-bin figures are the negative-weight/purity ratio over the
    populated cross-section bins.}
  \label{tab:negweight}
  \begin{tabular}{lcc}
    \hline\hline
    & binned purity & negative-weight \\
    \hline
    total $\sigma$ (cm$^2$/nucleon) & $\nwSigPur$ & $\nwSigNeg$ \\
    total ratio (neg/pur)           & \multicolumn{2}{c}{$\nwRatioTot$ (\SI{\nwPctTot}{\percent})} \\
    median per-bin ratio            & \multicolumn{2}{c}{$\nwMedianBin$} \\
    per-bin RMS deviation           & \multicolumn{2}{c}{\SI{\nwRmsBin}{\percent}} \\
    largest single-bin deviation    & \multicolumn{2}{c}{\SI{\nwWorstBin}{\percent} (low $p_T$, high background)} \\
    \hline\hline
  \end{tabular}
\end{table}

\begin{figure}[t]
  \centering
  \includegraphics[width=0.62\textwidth]{negweight_ratio_2d}
  \caption{Negative-weight/binned-purity cross-section ratio in populated
    bins of the 2D ME~FHC measurement, with common estimator and seed.
    The median is $\nwMedianBin$ (near-white). Only the low-$p_T$,
    high-background edge bin departs by more than a few percent. Background
    is largest there, and the $\max(0,\cdot)$/structural floor of
    \S\ref{sec:negweight-estimator} governs the result.}
  \label{fig:negweight-ratio}
\end{figure}

\subsection{Propagation through systematics and statistics}
\label{sec:negweight-uq}

Replacing binned subtraction requires reproducing its \emph{uncertainty}
budget as well as the central value. The full 2D ME~FHC covariance was rebuilt
end to end with negative-weight injection and compared to binned purity at
identical settings. The 187-universe systematic ensemble varied genuine
background consistently within each universe instead of freezing it at nominal;
a statistical bootstrap used matched seeds. Agreement is at the level predicted
by the $(D-B)/S$ identity. Overall systematic magnitude has
$\sqrt{\mathrm{tr}}$ ratio $\nwSystRatio$ (negative-weight/purity), and the
statistical ratio is $\nwStatRatio$.
Residuals of \SI{\nwSystResid}{\percent} and \SI{\nwStatResid}{\percent}
% --- each defined from its parent ratio rather than written
% out, so that a later revision of either ratio propagates to it ---
are consistent with ``to within \SI{2}{\percent}'' in \S\ref{sec:method}.
They are bootstrap and classifier-seed sampling noise, with no method bias.
Negative-weight injection reproduces the central cross section \emph{per bin}
(median $\nwMedianBin$, RMS \SI{\nwRmsBin}{\percent}, worst bin
\SI{\nwWorstBin}{\percent}) and covariance \emph{in aggregate}.
The covariance comparison gives two trace ratios; per-bin, correlation and eigenmode
equivalence are not established.

The 2D central-value, closed-form recovery and full covariance comparisons,
together with Stay-Positive robustness on larger extended-fiducial backgrounds,
support negative-weight injection as an unbinned alternative at the 2D scale.
Binned purity remains the default for this note's reproduction measurement.
Negative-weight injection is the validated unbinned cross-check and
\emph{designated} extended-fiducial baseline (\S\ref{sec:fps}), where
binning-independence and the continuous floor are essential.

\subsection{Which footing each production chain stands on}
\label{sec:negweight-footing}

Validation against one another does not make the two subtraction choices
arbitrary. The production chains choose differently.

\paragraph{The standard \(N\)-dimensional chain uses purity.}
The entire standard five-dimensional chain uses binned purity
(\S\ref{sec:method-algo}): the central unfold, $169$ vertical-band unfolds,
$18$ detector unfolds and $10$ active-lateral endpoint unfolds.
Its footing is positively identified, with different evidence for existing
unfolds and current production.

Existing endpoint unfolds from 2026-07 are identified by elimination.
Neither possible launcher passes \texttt{-{}-bkg-mode}; the driver defaults
to \texttt{purity} and announces its mode only on negative-weight branches.
The lateral logs contain the purity branch's signature line and no mode
announcement. The flag and announcements predate the files, making that absence
informative without a version gap. Current production pins the footing in a
configuration object that rejects every other mode
(App.~\ref{app:provenance}, \texttt{bkg-mode-pins}). The canonical launcher
passes \texttt{-{}-bkg-mode} explicitly from that pinned configuration
rather than inheriting the default, and the value is stamped into every endpoint
receipt as it is written. Asserting the previous default changes
provenance without changing physics. Its effect is that
a future component's footing is provable from the component instead of from the
absence of a flag.

\textbf{This is a recorded choice, not a silent default} (App.~\ref{app:provenance}, \texttt{gbdt-closeout} and
\texttt{oi6-oi8-oi126}). Two grounds carried it. First, mixing footings inside one covariance is worse than
either consistent choice: dropping a negative-weight-footed lateral block into a
purity-footed central and vertical set is a footing mismatch, not an upgrade. Second, switching required rerunning every
unfold above and would have invalidated the flux-corrected (J28) covariance
adopted immediately beforehand. That consequence decided the choice; cost did
not. The
decision carries an explicit obligation to revisit the footing before
submission. It is recorded as such, and nothing in this subsection discharges it.

\paragraph{What the quoted extended-fiducial number used.}
The $\sigma_{\mathrm{ext}}=\SI{\fpsSigma}{cm^2/nucleon}$ in \S\ref{sec:fps}
was not produced with negative-weight injection (App.~\ref{app:history}).
It comes from the FPS ME~FHC battery of 2026-06-10, before the negative-weight
driver. Its ten endpoint unfolds used the driver's purity default and are
recorded as purity controls, not publication inputs (App.~\ref{app:provenance},
\texttt{fps-battery}, \texttt{negweight-audit}). Negweight-refined FPS production
has not run. The ledger-verified central value from that gated battery remains
\SI{\fpsSigma}{cm^2/nucleon}. Negative-weight injection is the intended
extended-fiducial baseline; it did not produce the quoted number.

\paragraph{The published pair has different footings.}
The extended-fiducial lane designates negative-weight injection with Stay-Positive
refinement as its publication footing. Code pins it
(App.~\ref{app:provenance}, \texttt{bkg-mode-pins}) and a gate fails closed
if any endpoint's recorded mode differs or is missing. What the number in
\S\ref{sec:fps} actually used is the separate matter stated above;
the standard chain uses purity throughout. We claim internally consistent
footings for each side of the measurement and matched-control pair.
Uniformity across both is not claimed.

\subsection{Measured size of the footing choice and its scope}

\paragraph{Measured size.}
At full statistics with the adopted estimator on 2D ME~FHC, the covariance
difference is one to two percent. The $187$-universe systematic
$\sqrt{\mathrm{Tr}}$ ratio is $\nwSystRatio$, from
\SI{2.9828e-39}{cm^2/nucleon} negative-weight and
\SI{3.0242e-39}{cm^2/nucleon} purity. Matched-seed statistical covariance has
ratio $\nwStatRatio$, from \SI{1.7260e-40}{cm^2/nucleon} and
\SI{1.7576e-40}{cm^2/nucleon}, respectively. Real-data totals agree to
\SI{\nwPctTot}{\percent}, with per-bin median $\nwMedianBin$ and
\SI{\nwRmsBin}{\percent} RMS. These repeat the comparisons in
\S\ref{sec:negweight-2d} and \S\ref{sec:negweight-uq} with explicit operands
for checking the ratios.

\paragraph{Displacement across universes.}
The displacement argument rests on a \emph{measured} input; it assumes no
identity of continuous estimands. The constructions share bin masses in the
binned limit (\S\ref{sec:negweight-binned}) but differ by within-bin shape
terms (\S\ref{sec:negweight-estimator}). Switching footing produces a measured
net per-universe displacement of order \SI{0.1}{\percent}, including residual
and shape effects, at this measurement's sub-percent genuine-background scale.
Systematic covariance measures spread across universes about their own mean. Similar
small shifts of each universe leave that spread nearly unchanged, explaining
the $187$-universe ratio $\nwSystRatio$, far from the purity correction's
${\sim}\SI{3}{\percent}$. This argument connects a handful-of-universes
comparison to the $187$-universe covariance; the full $187$-universe rebuild
in \S\ref{sec:negweight-uq} tests it at 2D and passes.

\paragraph{Scope.}
Measured agreement concerns these 2D systematic and statistical covariances at
the stated configurations. It is not a theorem that the two
footings give the same covariance, because the underlying estimands are not
pointwise identical (App.~\ref{app:history}).
Covariance equivalence for the continuous estimator or in five dimensions is not established.
The five-dimensional evidence is limited to the following two-universe spot check.

\paragraph{Limits of the five-dimensional evidence.}
There is no full five-dimensional $187$-universe both-footing comparison at the
publication configuration ($5$ iterations, \texttt{lgbm}). The spot check
uses one iteration and \texttt{hist} on the on-disk five-dimensional universe
file. Under both footings \texttt{2p2h:0} gives $3.014\times10^{-38}$ with
per-bin median $1.0005$; \texttt{MaCCQE:0} differs by \SI{-0.03}{\percent}
with median $0.9990$. Together with the displacement argument and full-statistics
2D result, this supports quoting the standard chain on purity and stating its
footing. It does not support calling footing irrelevant in five dimensions;
that stronger claim requires the full comparison above. Raw negative weights
destabilize \texttt{exact} (\S\ref{sec:negweight-estimator}), while the standard
chain uses \texttt{lgbm}. Neither footing risks that pathology here, making
the choice safe without establishing that it is immaterial.

\subsection{Five-dimensional pseudo-experiments on both footings}

The five-axis validation (\S\ref{sec:val-fiveaxis})
reports end-to-end tests of both footings: purity bias in the highest-$W$ column
at nominal truth, its reduction with \texttt{negweight-refined}, the data
difference, and the refinement-capacity candidate, which is not adopted.

\paragraph{Purity bias.}
Purity bias is 8--12 statistical standard deviations; signal-only experiments
close. Subtraction is exact per reconstructed bin, but the widest bin
($W\ge\SI{3}{GeV}$, background \SI{13}{\percent} of signal) contains background
shape resolved by the unbinned classifier and missed by bin-uniform purity.
The within-bin shape term of \S\ref{sec:negweight-estimator} is therefore
measurable. It is direct evidence for the obligation in \S\ref{sec:negweight-footing}
to revisit footing before submission. Comparing five-dimensional purity with
\texttt{negweight-refined} on data is the cheapest estimate of its size on data.
% Source: docs/orchestration/OUTCOME-20260925-s5c-purity-background-bias-at-high-W.md; KNOWN_ISSUES 75.
% Heading date (2026-09-25) and the numbers repeated in sec_validation.tex removed (lane 3 App. B #7):
% highest-$W$ column \SIrange{3.5}{4.1}{\percent} low, highest-$\Eavail$, highest-$W$ cell \SI{1.7}{\percent} high.

\paragraph{Refined footing.}
Refinement used the unfold's deterministic GBDT configuration, independent
injection and subtraction halves of background simulation, and bootstrap
replicas that individually resampled every pseudo-event and refit refinement.
It preserves the signed total to $3\times10^{-6}$ and clips
\SIrange{0.01}{0.05}{\percent} of events. Corner bias falls from
\SI{1.66}{\percent} to \SI{0.24}{\percent}, but residuals in some joint cells
fail the predeclared nominal closure criterion. The data difference's sign
pattern matches simulated purity bias at an order of magnitude smaller size;
the difference is method sensitivity, not a measured bias. With either footing,
the estimator does not follow generator-sized $\Eavail$ shape changes.
No coverage test was run.
The footing obligation therefore remains open; reducing one bias does not itself qualify
an uncertainty.
% Source: docs/orchestration/OUTCOME-20260925-s5n-stage1-development-fail.md; VL151, VL152.
% Heading date (2026-09-25) and numbers repeated in sec_validation.tex removed (lane 3 App. B #7):
% highest-$W$ bias falls from \SIrange{-3.5}{-4.1}{\percent} to \SIrange{0.11}{0.24}{\percent};
% residual of up to about \SI{1}{\percent}; on the data the two footings
% differ by a median of \SI{0.36}{\percent} and at most \SI{1.4}{\percent}.

\paragraph{Refinement capacity.}
Stay-Positive classifier underfitting causes the refined footing's residual
nominal bias. At the unfold's GBDT setting (100 trees, 8 leaves), the
positive-weight probability is underfit, leaving up to \SI{4}{\percent}
reco-level subtraction error per $(\Eavail,W)$ reco cell. With 400 trees and 31 leaves,
error falls to \SI{1.7}{\percent}, paired truth-level bias falls by more than
\SI{80}{\percent} in every affected cell, and nominal background-inclusive
closure passes on 20 development and 40 fresh pseudo-experiments.
Splitting background simulation into independent injection and subtraction
halves is not the cause.
% Source: docs/orchestration/DIAGNOSIS-20260926-s5e-oi192-estimator.md §5;
% docs/orchestration/OUTCOME-20260926-s5e-oi192-diagnosis-and-candidate.md; VL153, VL154.
% Heading date (2026-09-26) removed. Moved to the pointer above: "The higher-capacity refinement
% is not adopted: the estimator it feeds failed a stability test on the data (\S\ref{sec:validation})."

\paragraph{The footing was chosen, not discharged; the adoption gate is separate.}
Nothing here adopts any five-dimensional covariance. The selection-complete
active-lateral component is built, with five bands and ten endpoint unfolds
carrying distinct digests bound in the adopted product's receipt. One five-axis
covariance is adopted under a digest-bound exception
(\S\ref{sec:eavailw-adopted}). Neither fact follows from or is needed by this
appendix. Footing choice and adoption are independent; the choice removes one
blocking decision while leaving every other component precondition intact
(App.~\ref{app:provenance}, \texttt{p4-verifier}).
